\section{向量执行：一条指令怎样推进多个数据元素}

超标量处理器在同一周期发射多条彼此独立的标量指令；\term{单指令多数据}{single instruction, multiple data, SIMD} 则让一条指令对多个数据元素执行同一种操作。二者可以同时存在：前端把向量指令译成一个或多个微操作，乱序核心仍负责重命名、唤醒和提交，而向量执行单元在多个 lane 上并行处理元素。Intel 的 SSE、AVX 系列和 RISC-V V 扩展采用的寄存器宽度与控制方式不同，但都必须解决元素宽度、有效 lane、掩码、访存和精确异常这些共同问题。\cite{intel-sdm,riscv-vector}

\topic{向量状态不仅是一组更宽的寄存器}
一条向量加法的输入状态至少包括两个源向量、目标向量、元素宽度以及有效元素集合。固定宽度 SIMD 通常由指令编码决定寄存器宽度和元素类型；向量长度可变的 ISA 还保存向量长度和元素配置状态。掩码决定哪些 lane 可以写回，未激活 lane 必须按 ISA 规定保持旧值、写零或处于无关状态，不能由具体实现任意选择。

数据路径从向量寄存器文件或重命名后的物理向量寄存器开始。源操作数经过 lane 划分进入加法、乘法、置换或转换单元；结果先进入完成缓冲或物理目标寄存器，最终仍由 ROB 的按序退休边界变成架构可见状态。若 256-bit 数据路径处理 32-bit 元素，一次可容纳 8 个元素；但“容纳 8 个”不等于“每周期完成 8 个”。实际吞吐还取决于执行端口数量、指令延迟、单元是否分拍复用以及目标寄存器的数据相关。

\topic{向量 load/store 仍受 Cache line、页和 LSQ 约束}
连续向量加载可以在一个或多个 Cache line 内形成若干访问。地址生成单元先计算起始地址和跨度，LSQ 检查与更早 store 的顺序，再由 Cache 端口拆分跨行请求。若向量横跨页边界，两个页面必须分别完成 TLB 查找和权限检查。前半段命中、后半段缺页时，处理器不能把一个不符合 ISA 异常语义的“半更新结果”提交给软件。

聚集/散布（gather/scatter）进一步把每个 lane 的索引变成独立地址。它能表达不规则访问，却可能产生多个 TLB 查找、多个 Cache miss 和多次 bank 冲突。此时向量算术单元并不是瓶颈，地址生成、MSHR、TLB 和内存带宽才决定完成时间。性能分析必须区分“每条指令处理的元素数”和“每周期真正返回的元素数”。

\topic{掩码与精确异常共同定义提交边界}
设一条掩码向量除法有 8 个 lane，其中 5 个 lane 有效。无效 lane 不应触发除零异常；有效 lane 中若第 4 个访问了无效页面，处理器必须按照 ISA 定义保存可重启状态或在退休点报告异常。实现可以让较早 lane 已经完成计算，但在异常指令退休以前，这些结果仍不能越过架构提交边界。恢复时还要知道哪些元素已经完成、哪些需要重做，因此向量异常状态本身也是处理器与操作系统之间的接口。

\topic{上下文切换要保存实际可见的向量状态}
操作系统不能仅保存通用寄存器。线程一旦使用向量扩展，其向量寄存器、掩码和控制状态也属于架构上下文。更宽的状态会增加保存、恢复和迁移成本；实现可以采用延迟保存等策略，但必须确保另一线程无法观察前一线程遗留的架构状态。这里再次出现本章的共同边界：微架构可以延后搬运，线程重新运行前的可见状态却必须完整恢复。

\section{SMT：多个硬件线程怎样共享一个乱序核心}

\term{同时多线程}{simultaneous multithreading, SMT} 在一个物理核心内维护两组或更多架构线程状态，并允许它们的微操作共同竞争前端、执行端口和存储层次。每个逻辑线程至少具有独立 PC、架构寄存器映射和异常状态；分支预测器、取指带宽、译码器、物理寄存器、发射队列、执行单元、Cache 与 TLB 则可能按实现选择共享、分区或动态借用。SMT 的目标不是让一个线程凭空获得更多执行单元，而是让另一个线程利用第一个线程因依赖或 Cache miss 留下的空槽。\cite{arm-smt}

\topic{逐周期共享从取指选择开始}
假设线程 A 因 L2 miss 暂时缺少可发射微操作，线程 B 的操作数已经就绪。取指策略可以暂时提高 B 的份额，重命名阶段为 B 分配物理寄存器和 ROB 项，发射器再把 B 的微操作送入空闲端口。A 的数据返回后，两线程同时具有就绪工作，仲裁器必须限制单个线程长期占满 ROB、LSQ 或提交带宽，否则另一线程可能无法前进。

因此 SMT 的控制状态至少包含每线程取指优先级、各结构占用计数和公平性阈值。一次微操作退休时，只更新所属线程的架构映射；一个线程发生分支误预测或异常时，也只能清除该线程中更年轻的工作，不能破坏另一线程已经建立的状态。

\topic{吞吐增加与单线程性能下降可以同时发生}
若单线程只能利用 60\% 的执行槽，加入第二线程可能把核心总利用率提高到 85\%。但两个线程还会争用 L1、TLB、分支预测表、内存带宽和功耗预算，所以线程 A 的单独完成时间可能变长。评价 SMT 必须同时报告总吞吐、每线程尾延迟、公平性以及共享资源 miss，而不能把“两个逻辑处理器均显示忙碌”直接解释为性能翻倍。

\topic{共享状态也是隔离边界}
同一物理核心上的线程可能通过 Cache、预测器、端口争用和执行时延相互影响。调度器在安排不同安全域时，需要结合平台提供的预测控制、Cache 隔离和 SMT 策略判断风险。关闭 SMT 可以消除一部分同核共享，但不能自动消除跨核心共享 Cache、内存控制器或互连产生的全部侧信道。具体安全边界将在“推测执行留下的微架构痕迹”一节继续展开。
